\section{Introduction}
\label{sec:intro}

Runtimes that mediate untrusted code have acquired demanding new
customers. Confidential-compute platforms expect every cross-component
call to be brokered with an authorisation story a reviewer can
actually audit~\cite{tdx2024,sevsnp,arm-cca}. The AI agents that now
drive browsers, shells, and developer toolchains in production route
every tool call through some flavour of \texttt{bubblewrap},
container, or microVM~\cite{anthropic-cu,openai-agents,devin}---and
none of those expose a capability layer above the syscall boundary.
Regulators, meanwhile, have started writing the requirement down:
NIST AI~600-1 and Article~15 of the EU AI Act both push toward
\textit{verifiable mediation} of agent
actions~\cite{nist-ai-600-1,euaiact15}, which in practice means
authorisation that is cheap enough to sit on every message and
unforgeable enough to mean something.

The systems that come closest today divide along two axes
(\Cref{tab:landscape}): \textbf{kernel-resident vs.\ userspace}, and
\textbf{capability-mediated vs.\ not}. seL4~\cite{klein2009sel4} owns
the kernel-capability corner; its assurances are the strongest
available, but it is a clean-slate kernel and cannot be layered onto
Linux. gVisor~\cite{gvisor} and Firecracker~\cite{firecracker} hold
the userspace-mediation corner without any capability model at
all---isolation comes from syscall interception or microVM
separation. CAP-VMs~\cite{capvms} pairs capabilities with fast IPC in
userspace, but only on CHERI/Morello hardware. The corner that stays
empty---the one the agent-runtime workload keeps pushing toward---is
\textit{userspace, capability-mediated, stock Linux, ordinary
hardware}.

\sysnamefull is our attempt to occupy that corner. It is a C++
runtime that gates every zero-copy memfd IPC channel through
unforgeable 128-bit \captokens with epoch-based cascading revocation,
and it delivers validated messages in under a microsecond without
touching the kernel.

\para{The empirical question}
A capability check that costs hundreds of nanoseconds per message
kills the idea on arrival; nobody adopts a mediated fastpath that
throws away the latency the fastpath existed for. So the question
this paper answers is not \emph{can we gate every message?}---of
course we can---but \emph{can the gate ride the fastpath without
giving up the floor that makes zero-copy IPC worth having?} We answer
it by measurement: five reference transports---pipe(2), socketpair,
io\_uring, Aeron, eRPC---run head-to-head in the same TSC domain on
the same machine.

\para{Why this is not obvious}
Two design decisions are easy to get wrong. The natural
implementation of validate is a hash-table lookup, which costs an
indirection and usually a cache miss---on a sub-microsecond budget,
that is most of the budget. And the natural threat model admits
\textit{forged token identifiers}, which seems to rule out the one
obvious shortcut: trusting the token to say where it lives in the
pool.

Our way out is to trust the token \emph{partially}. The high 32 bits
of the 128-bit identifier carry the pool slot index; the low 96 bits
are RDRAND output. The slot index is a hint, never a credential---a
constant-time compare over the full 128 bits still has to pass before
the token is accepted. Forgery resistance stays at $2^{96}$; the
lookup drops to a shift and a load.

\para{Contributions}
\begin{enumerate}
\itemsep0pt
\item The first software-only realisation on stock Linux of a
Redell\,/\,EROS\,/\,Cornucopia\,Reloaded~\cite{redell74,
shapiro1999eros,cornucopia-reloaded}-class revocable-capability
regime, gating a wait-free MPSC IPC fastpath in shared mmap memory
(\Cref{sec:design}, \Cref{sec:impl}).
\item The slot-indexed validate path that makes the gate affordable
without shrinking the $2^{96}$ forgery space
(\Cref{sec:design:fastpath}).
\item A reproducible same-TSC-domain comparison against pipe(2),
socketpair, io\_uring, Aeron, and eRPC, with throughput, MPSC
scaling, revocation-latency, and perf-counter measurements, all
driven by one shell command (\Cref{sec:eval}).
\item A property-based oracle test, a CBMC bounded-model proof of
permission monotonicity, and an MPSC-contention revocation test that
together back every numerical claim we make
(\Cref{sec:eval:correctness}).
\end{enumerate}

\Cref{sec:eval} reports a per-message gate cost of
\gateOverheadTail\,ns at p99 across payload sizes from 64\,B to
16\,KiB. That is \gateOverheadPercent\% added to an ungated round
trip of \rawAstraRTT\,ns---a large fraction of a very small
number---and the mediated path that results, at
\gatedAstraRTT\,ns, is still faster than every transport in our
comparison set. Mediation, in other words, does not cost you the
fastpath. The validate itself runs in \validateTail\,ns at p99
\textit{regardless of pool occupancy} (\Cref{fig:figure-2}), which is
the $O(1)$ claim made concrete.

\para{What we do not claim}
$O(1)$ revocation is not our algorithm. Redell~\cite{redell74} and
EROS~\cite{shapiro1999eros} established the pattern half a century
ago, and Cornucopia\,Reloaded~\cite{cornucopia-reloaded} is the
modern peer on CHERI. Our contribution is the first software-only
realisation on stock Linux, with measured numbers under contention.
The same goes for memfd zero-copy IPC as a substrate: Aeron and
io\_uring ship the same template, and have for years. We reuse the
substrate; what is new is the gate that sits on it.

\para{Roadmap}
\Cref{sec:threat} states the threat model. \Cref{sec:bg} positions
the work against capability OSes and modern IPC primitives.
\Cref{sec:design} develops the slot-indexed gate; \Cref{sec:impl}
covers the C++ implementation, the M-21 constant-time NASM
primitives, and the CBMC proof. \Cref{sec:eval} presents the
evaluation, \Cref{sec:disc} the limitations---including our choice of
C++---and \Cref{sec:rw} the related work in detail.
\Cref{sec:conclusion} concludes.
